\section{Extracting a tree packing with small resistance}
\label{sec:extraction}

Throughout this section, \(H\) is an \(r\)-tree-packable graph on
\(n\ge2\) vertices. Set
\begin{equation}\label{eq:parameters}
 u=r^{-1/16},\qquad q=\lfloor(1-u)r\rfloor.
\end{equation}
We take \(r\) above an absolute threshold, so in particular \(q>1\).
The threshold may be increased finitely many times below.

\begin{theorem}[Shortcut extraction]\label{thm:extraction}
For all sufficiently large integers \(r\), every \(r\)-tree-packable
graph \(H\) admits \(X\succ0\) such that
\begin{equation}\label{eq:shortcut-X}
 L_H\preceq X,\qquad
 \mathbf1_S^{\mathsf T}X\mathbf1_S\le3|\delta_H(S)|
 \quad(\varnothing\ne S\subsetneq V),
\end{equation}
and the edges with \(R_X(e)\le4r^{-1/2}\) contain \(q\)
edge-disjoint spanning trees. The lower threshold on \(r\) is
independent of \(H\) and \(n\).
\end{theorem}

For a nontrivial maximal packing partition,
Lemma~\ref{lem:packing-parts}(iii) limits the number of crossing edges
in the threshold graph. We first construct
a response that makes most crossing edges have small resistance.
Averaging responses at nearby matrices and taking a fixed-point limit transfers
this estimate to the matrix defining the partition, forcing that
partition to be trivial.

\subsection{A compact set of matrices}

Let
\begin{equation}\label{eq:compact-domain}
 \mathcal K=
 \left\{X=X^{\mathsf T}:
 X\succeq L_H+n^{-1}I,\quad
 \mathbf1_S^{\mathsf T}X\mathbf1_S\le3|\delta_H(S)|
 \text{ for }\varnothing\ne S\subsetneq V\right\}.
\end{equation}
This is a nonempty compact convex subset of the real symmetric
matrices. Indeed, \(H\) is connected, and
\(|S|/n<1\le|\delta_H(S)|\), so \(L_H+n^{-1}I\) belongs to
\(\mathcal K\). Closedness and convexity follow from the defining
inequalities. The singleton cuts give
\(X_{vv}\le3\deg_H(v)\), and positive semidefiniteness gives
\(|X_{vw}|^2\le X_{vv}X_{ww}\); hence the domain is bounded.
Every matrix in \(\mathcal K\) is bounded below by \(n^{-1}I\).

Given any locally \(q\)-connected hierarchy in \(H\), mark exactly
the nonroot nodes satisfying
\[
 |O(A)|\ge q^{-1}|P(A)|.
\]
For all sufficiently large absolute \(r\), we have \(q\ge r/2\)
and \(q\ge q_H\). Theorem~\ref{thm:hierarchy-power} gives a shortcut
matrix \(D\) whose average resistance on every marked \(O(A)\) is
at most
\[
 q^{-31/40}\le2^{31/40}r^{-31/40}\le r^{-3/4},
\]
where the final inequality holds for \(r\ge2^{31}\). Put
\begin{equation}\label{eq:hierarchy-response}
 Y=L_H+n^{-1}I+D.
\end{equation}
Then \(Y\in\mathcal K\), since on every cut its quadratic form is
at most \(2|\delta_H(S)|+|S|/n\le3|\delta_H(S)|\).
Lemma~\ref{lem:resistance-order} also gives
\begin{equation}\label{eq:response-average}
 \frac1{|O(A)|}\sum_{e\in O(A)}R_Y(e)\le r^{-3/4}
 \qquad\text{for every marked }A.
\end{equation}

\subsection{Packing partitions and a simultaneous response}

Fix \(X\in\mathcal K\). For \(s>0\), let
\[
 F_s(X)=\bigl(V,\{e\in E(H):R_X(e)\le s\}\bigr),
 \qquad
 \mathcal P_s(X)=\mathcal P(F_s(X)),
\]
where \(\mathcal P(F)\) is the maximal packing partition of
Lemma~\ref{lem:packing-parts}. Define
\[
 B_s(X)=E_H(\mathcal P_s(X)),\qquad M_s(X)=|B_s(X)|.
\]
The crossing sets contain edges of \(H\), including those absent
from \(F_s(X)\).
As \(s\) increases, the partitions coarsen and their crossing sets
decrease. Only finitely many partitions occur.
Each part induces a \(q\)-edge-connected subgraph of \(H\),
because it contains \(q\) packed trees in \(F_s(X)\).

\begin{lemma}[Simultaneous response]\label{lem:response}
There is an absolute constant \(C_0\), which may be taken to be \(40\),
such that for every \(X\in\mathcal K\) there exists \(Y\in\mathcal K\)
with
\begin{equation}\label{eq:simultaneous-response}
 \bigl|\{e\in B_s(X):R_Y(e)>r^{-1/2}\}\bigr|
 \le C_0r^{-1/4}M_s(X)
 \qquad\text{for every }s>0.
\end{equation}
\end{lemma}
\begin{proof}
Temporarily suppress \(X\) from the notation.
In each occupied bin \(2^a\le M_s<2^{a+1}\), \(a\ge0\),
choose the finest partition that occurs in that bin.
List the chosen partitions in order from fine to coarse as
\(\mathcal P_0,\ldots,\mathcal P_{\ell-1}\), with crossing sets
\(B_i\) and sizes \(M_i\). Append
\(\mathcal P_\ell=\{V\}\), \(B_\ell=\varnothing\), \(M_\ell=0\).

The distinct parts of these partitions, together with all
singletons, form a laminar family containing \(V\).
Give each proper set its smallest proper containing set as parent.
The children of every nonsingleton partition that set, so this is
a hierarchy. Every induced graph at a node is \(q\)-edge-connected:
this was proved for the partition parts, the root is
\(r\)-tree-packable, and singleton conditions are vacuous.
Take the response \(Y\) from \eqref{eq:hierarchy-response} for this
hierarchy.

Consider a transition \(\mathcal P_j\) to \(\mathcal P_{j+1}\).
If a part \(A\in\mathcal P_j\) changes, the part of
\(\mathcal P_{j+1}\) containing it is its immediate parent in the
hierarchy. Indeed, earlier selected partitions are finer and later
ones are coarser, so none supplies a distinct set strictly between
these two. Adding singleton leaves creates no such set.

An edge of \(B_j\setminus B_{j+1}\) has endpoints in two changing
parts of \(\mathcal P_j\) with the same parent. It belongs to the
\(O(A)\) set of each of those parts; see Figure~\ref{fig:hierarchy}.
The total full boundary size of all changing parts is at most
\(2M_j\), since they belong to the single partition \(\mathcal P_j\).
For the unmarked changing parts,
\[
 \sum_A |O(A)|\le \frac1q\sum_A|P(A)|\le\frac{2M_j}{q}.
\]
For the marked changing parts, \eqref{eq:response-average} gives
\[
 \bigl|\{e\in O(A):R_Y(e)>r^{-1/2}\}\bigr|
 \le r^{-1/4}|O(A)|.
\]
Their total bad-edge count is therefore at most \(2M_jr^{-1/4}\).
These counts may cover an edge more than once, which only enlarges
the upper bound. Consequently at most
\[
 2(r^{-1/4}+1/q)M_j
\]
of the disappearing edges are bad.

There is only one representative per occupied dyadic bin, so its
bin exponents strictly decrease. For every \(i<\ell\),
\begin{equation}\label{eq:dyadic-tail}
 \sum_{j=i}^{\ell-1}M_j<4M_i.
\end{equation}
For example, if \(2^a\le M_i<2^{a+1}\), the successive terms are
bounded by \(2^{a+1},2^a,2^{a-1},\ldots\).
Every edge of \(B_i\) disappears at exactly one later transition.
Its bad-edge count is thus at most
\[
 8(r^{-1/4}+1/q)M_i.
\]
For an arbitrary \(s\) with \(M_s>0\), the chosen representative
in its bin is finer, so \(B_s\subseteq B_i\) and \(M_i<2M_s\).
The bad-edge count in \(B_s\) is at most
\[
 16(r^{-1/4}+1/q)M_s\le40r^{-1/4}M_s
\]
for all sufficiently large \(r\). If \(M_s=0\), the assertion is
immediate. This proves \eqref{eq:simultaneous-response}
simultaneously for every threshold.
\end{proof}

\begin{figure}[tbp]
  \centering
  \begingroup
  \definecolor{hierarchyblue}{RGB}{37,87,128}
  \begin{tikzpicture}[
    x=1cm,y=1cm,
    font=\small,
    coarse/.style={draw=black,line width=.8pt,rounded corners=5pt},
    fine/.style={draw=black!65,dashed,line width=.65pt,rounded corners=4pt},
    vertex/.style={circle,fill=black,inner sep=1.5pt},
    levelarrow/.style={-{Stealth[length=1.8mm]},line width=.55pt},
    every node/.style={inner sep=2pt}
  ]
    \node at (6.35,6.25) {Selected threshold partitions: fine to coarse};
    \node (pzero) at (1.10,5.78) {$\mathcal P_0$};
    \node (dotsleft) at (2.85,5.78) {$\cdots$};
    \node (pi) at (4.60,5.78) {$\mathcal P_i$};
    \node (pinext) at (7.00,5.78) {$\mathcal P_{i+1}$};
    \node (dotsright) at (9.00,5.78) {$\cdots$};
    \node (root) at (11.05,5.78) {$\{V\}$};
    \draw[levelarrow] (pzero) -- (dotsleft);
    \draw[levelarrow] (dotsleft) -- (pi);
    \draw[levelarrow] (pi) -- (pinext);
    \draw[levelarrow] (pinext) -- (dotsright);
    \draw[levelarrow] (dotsright) -- (root);

    \draw[fine] (1.20,5.12) -- (2.00,5.12);
    \node[anchor=west] at (2.10,5.12) {parts of $\mathcal P_i$};
    \draw[coarse] (7.15,5.12) -- (7.95,5.12);
    \node[anchor=west] at (8.05,5.12) {parts of $\mathcal P_{i+1}$};

    \draw[coarse] (.25,1.40) rectangle (6.15,4.65);
    \draw[coarse] (6.55,1.40) rectangle (12.45,4.65);
    \node at (3.20,4.33)
      {$C=A^*=(A')^*$};
    \node at (9.50,4.33) {$D$};

    \draw[fine] (.65,1.85) rectangle (2.95,3.98);
    \draw[fine] (3.45,1.85) rectangle (5.75,3.98);
    \draw[fine] (6.95,1.85) rectangle (9.25,3.98);
    \draw[fine] (9.75,1.85) rectangle (12.05,3.98);
    \node at (1.80,3.65) {$A$};
    \node at (4.60,3.65) {$A'$};
    \node at (8.10,3.65) {$D_1$};
    \node at (10.90,3.65) {$D_2$};

    \node[vertex] at (1.20,2.48) {};
    \node[vertex] at (1.45,3.10) {};
    \node[vertex] (u) at (2.25,2.80) {};
    \node[vertex] (v) at (4.15,2.80) {};
    \node[vertex] at (4.95,3.10) {};
    \node[vertex] (w) at (5.10,2.35) {};
    \node[vertex] (z) at (7.60,2.35) {};
    \node[vertex] at (7.70,3.10) {};
    \node[vertex] at (8.60,2.75) {};
    \node[vertex] at (10.35,2.50) {};
    \node[vertex] at (10.60,3.10) {};
    \node[vertex] at (11.40,2.80) {};

    \draw[hierarchyblue,line width=1.3pt] (u) -- (v)
      node[midway,above=3pt,text=hierarchyblue] {$e$};
    \draw[black,line width=.7pt] (w) -- (z)
      node[midway,above=3pt] {$f$};

    \node[text=hierarchyblue] at (6.35,.86)
      {$e\in B_i\setminus B_{i+1},\qquad e\in O(A)\cap O(A')$};
    \node at (6.35,.30)
      {$O(A)=E_H(A,C\setminus A),\qquad f\in B_{i+1}\subseteq B_i$};
  \end{tikzpicture}
  \endgroup
  \caption{At a selected transition $\mathcal P_i\to\mathcal P_{i+1}$,
    the parts $A,A'$ merge into their immediate parent $C$.
    The edge $e$ leaves the crossing set and belongs to both $O(A)$ and
    $O(A')$, whereas $f$ remains a crossing edge.
    Only representative vertices and edges are shown.}
  \label{fig:hierarchy}
\end{figure}


\subsection{A consistent matrix from a fixed-point limit}

\begin{proof}[Proof of Theorem~\ref{thm:extraction}]
The response supplied by Lemma~\ref{lem:response} need not vary
continuously with \(X\). We interpolate only finitely many responses
at a time.

For each positive integer \(p\), cover \(\mathcal K\) by finitely
many relative open balls of radius \(1/p\), with centers
\(C_{p,1},\ldots,C_{p,N_p}\in\mathcal K\). At each center choose a
response \(Y_{p,j}\). Define continuous weights
\[
 \phi_{p,j}(Z)=
 \frac{\max\{0,p^{-1}-\|Z-C_{p,j}\|\}}
 {\sum_{h=1}^{N_p}\max\{0,p^{-1}-\|Z-C_{p,h}\|\}}
 \qquad(Z\in\mathcal K).
\]
The denominator is positive by the open cover.
The map
\[
 \Phi_p(Z)=\sum_{j=1}^{N_p}\phi_{p,j}(Z)Y_{p,j}
\]
is a continuous self-map of the compact convex set \(\mathcal K\).
Theorem~\ref{thm:fixedpoint} gives
\(X_p=\Phi_p(X_p)\).
Pass to a subsequence along which \(X_p\to X\in\mathcal K\).

Choose
\[
 3r^{-1/2}<s<4r^{-1/2}
\]
distinct from every \(R_X(e)\), \(e\in E(H)\).
This is possible because the graph is finite.
If \(\phi_{p,j}(X_p)>0\), then
\[
 \|C_{p,j}-X\|<p^{-1}+\|X_p-X\|\longrightarrow0
\]
uniformly over these centers. The resistance of each edge is
continuous on \(\mathcal K\). Since all resistances at \(X\) avoid
\(s\), for all sufficiently large \(p\) every contributing center
has exactly the same threshold graph, packing partition, and
crossing set as \(X\).

Write \(B=B_s(X)\), \(M=|B|\), and
\(\delta=C_0r^{-1/4}\). If \(M=0\), connectedness of \(H\) implies
that its partition is \(\{V\}\), and the conclusion already follows.
Otherwise, for large \(p\), each contributing \(Y_{p,j}\) has at
most \(\delta M\) bad edges in this same set \(B\).
Summing their bad indicators with weights \(\phi_{p,j}(X_p)\)
shows that at most \(2\delta M\) edges receive bad weight greater
than one half.

For every other edge \(e\in B\), at least half the total weight is
on matrices satisfying
\[
 Y_{p,j}\succeq r^{1/2}b_eb_e^{\mathsf T},
\]
by Lemma~\ref{lem:resistance-order}. All remaining matrices are
positive semidefinite. Hence their convex combination satisfies
\[
 X_p\succeq\tfrac12r^{1/2}b_eb_e^{\mathsf T},
 \qquad R_{X_p}(e)\le2r^{-1/2}.
\]
Thus at least \(N=\lceil(1-2\delta)M\rceil\) edges satisfy this
closed inequality. The condition that at least \(N\) edges of the
finite set \(B\) obey it is closed: it is a finite union, over
\(N\)-element subsets of \(B\), of finite intersections of closed
conditions.
The same count therefore holds at the limit \(X\).
Since \(2r^{-1/2}<s\), these edges belong to \(F_s(X)\).

Suppose \(\mathcal P_s(X)\) has \(d>1\) parts.
The original \(r\) packed trees in \(H\) give \(M\ge r(d-1)\).
Choose the absolute threshold on \(r\) so that
\(2C_0r^{-1/4}<u\). Then
\[
 |E_{F_s(X)}(\mathcal P_s(X))|
 \ge(1-2C_0r^{-1/4})M
 >q(d-1),
\]
contradicting Lemma~\ref{lem:packing-parts}(iii).
The partition is therefore \(\{V\}\), so \(F_s(X)\) is
\(q\)-tree-packable. Every one of its edges has resistance at most
\(s<4r^{-1/2}\), and \(X\in\mathcal K\) gives
\eqref{eq:shortcut-X}.
\end{proof}

\begin{remark}\label{rem:uniform-extraction}
For a fixed graph, the finite-cover scale needed to stabilize the
chosen threshold can be arbitrarily small. This does not affect the
absolute cutoff on \(r\). The latter comes only from the hierarchy
estimate and the numerical inequalities in the proof.
The construction makes no claim about its computational cost.
\end{remark}
